\section{Case Studies}
\label{sec:case-studies}

Aggregate numbers (\S\ref{sec:evaluation}) show \emph{that} the
pipeline performs well; the five cases below show \emph{why}: concrete
instances of the validation discipline finding and resolving a real
problem, each illustrating a different class of failure a purely
numeric evaluation would not surface on its own.

\subsection{RoningLoader's Multi-Stage Chain}
\label{sec:case-roning}

RoningLoader's infection chain, established through manual reverse
engineering, is not one binary but four: a loader that decrypts and
drops a kernel-mode rootkit driver used to hide its own process and
files, a separate DLL whose sole function is enumerating and
terminating security products, and a remote-access client that
establishes the actual C2 channel. Each component's own capability is
real and independently significant; none of the three dropped
components appear as an import or string in the loader binary itself,
since they are decrypted into memory or dropped to disk only at
runtime.

Table~\ref{tab:results} already shows the numeric consequence:
scoring the loader in isolation reaches 0.68 recall against ground
truth written for the full chain, because roughly a third of the
techniques a human analyst correctly attributes to ``RoningLoader'' as
a family are only observable by analyzing the dropped components
directly. The resubmission loop (\S\ref{sec:resubmission}) exists
specifically because this is a structural property of multi-stage
malware, not a quirk of one sample, and closes 16 of those 32 recall
points by giving each component its own independent analysis pass
rather than attempting to infer its capability from the parent's own
evidence.

\subsection{The Installer Aggregation Bug}
\label{sec:case-installer}

Early in validation, an installer-packaged sample's static report
showed evidence consistent with only \emph{one} of its several bundled
components, despite the multi-pass extraction logic
(\S\ref{sec:static-rules}) having processed all of them. Tracing the
report generation code directly, rather than assuming the multi-pass
logic itself was at fault, found the actual defect one stage later:
\code{\_analyze\_installer}, the function aggregating per-component
results into the sample's single top-level report, was silently
\emph{substituting} one child component's analysis for another's under
a specific condition, rather than merging both.

The bug was confirmed, not just inferred from the symptom, via the
report's own \code{hash\_match} field, a value recorded for exactly
this kind of cross-check, showing the reported hash belonged to a
different extracted file than the one whose findings were actually
present in the report body. Fixing the aggregation logic to merge
findings from every extracted component, rather than the most recently
processed one, restored the missing evidence and directly improved
recall on the affected sample once re-evaluated.

The lesson generalized beyond this one fix: it motivated treating every
new pipeline capability, the resubmission loop chief among them, as
needing its own explicit corroboration step (lineage manifests, hash
cross-checks) rather than trusting that ``the code ran without error''
implies ``the code produced the right output.'' Finding and fixing this
class of bug before it could silently understate a sample's real
capability is itself a concrete demonstration of what the pipeline's
validation discipline is for.

\subsection{WhiteSnakeStealer's Ground-Truth Gap}
\label{sec:case-wsnake}

\S\ref{sec:fp-audit} summarizes the outcome; the process is worth
detailing once as a concrete example of what ``audited against the full
report text'' means in practice. WhiteSnakeStealer's pipeline output
included \technique{T1620} (Reflective Code Loading) with static
evidence, specific .NET BCL calls consistent with in-memory assembly
loading. The sample's own ATT\&CK Techniques table, however, had no
\technique{T1620} row, so the automated evaluation initially scored
this as a false positive.

Reading the report's narrative section, rather than trusting its
summary table alone, found an explicit paragraph describing the exact
reflective-loading mechanism the static rule had flagged: independent
analyst confirmation existed, it simply had not been transcribed into
the table the automated ground-truth extractor reads. This is a
documentation gap in the source report, not a pipeline error, and was
fixed by adding the missing row to the report itself (and regenerating
ground truth from it), rather than by adjusting the pipeline's
confidence or suppressing the finding. It is presented here
specifically because two structurally similar cases on RoningLoader
(\S\ref{sec:fp-audit}) were resolved the opposite way, left open,
precisely because this kind of independent textual confirmation was
what those two cases lacked. Together, the two outcomes show the
pipeline's own evidence being treated the same way regardless of
whether the discrepancy ultimately favored it.

\subsection{The XOR Key False-Positive Cascade}
\label{sec:case-xor}

Single-byte XOR brute-forcing (\code{config\_parser.py}) exhaustively
tries all 256 keys against candidate ciphertext regions and
regex-matches the result against a dotted-quad IP pattern, the
technique that recovered RoningLoader's genuine C2 IP
(\code{202.95.11.173}) with no prior key knowledge
(\S\ref{sec:limitations}, item~\ref{lim:key-ceiling}). Run without safeguards against a
second sample, WhiteSnakeStealer, the same routine produced 8
syntactically valid but entirely spurious dotted-quad matches from a
single wrong key (\code{0x2e}, ASCII \code{.}): a repeating,
null-padded integer table in the binary's compiled data happened to
decode to eight fake IP addresses under that one key, since null bytes
decode to literal \code{.} characters under it and a repeating
structure in the source data produces a repeating false match.

Two filters closed the gap, each justified by a distinct property a
genuine embedded value has and a decoding accident does not. First, a
\textbf{clean-boundary check}: the byte immediately preceding and
following a match must be \code{0x00} or the key byte itself, which a
real, null-terminated embedded string satisfies and an arbitrary byte
sequence decoded by an unrelated key generally does not. Second, a
\textbf{match-density cap} of two IP-shaped matches per key anywhere in
the file, on the principle that a genuine configuration value appears
once, or rarely a handful of times, while a cascade from decoding
unrelated structured data under the wrong key does not self-limit.
Verified against both directions afterward: RoningLoader's genuine C2
IP still produces exactly one match and survives the density cap;
WhiteSnakeStealer's eight spurious matches are now correctly discarded.

The lesson generalized beyond this one fix: it is what motivated the
explicit project policy (\S\ref{sec:limitations}, item~\ref{lim:tool-constraints}) that any new
exhaustive-search-based recovery capability must be validated against a
second sample, not just the one that motivated building it, before
being trusted. Had this capability been declared production-ready after
RoningLoader alone, the pipeline would have shipped with a
false-positive generator that happened not to have fired yet.

\subsection{Network IOC Contamination: Parsing Noise and Sandbox
  Background Traffic}
\label{sec:case-ioc-noise}

Auditing the pipeline's STIX/MISP-facing IOC output directly, rather
than only the ATT\&CK technique scores \S\ref{sec:evaluation} already
covers, surfaced two independent contamination sources neither
technique-level metric can see, since neither changes which techniques
are reported: fixing both left every sample's technique set, precision,
and recall bit-for-bit unchanged.

First, the static domain extractor's dotted-string regex was filtered
only by a hand-written denylist of namespace prefixes anchored to the
start of the match, so it mis-tokenized compiler metadata as domains:
\code{kernel32.dllfailed}, \code{asio.misc}, and, across RoningLoader's
26 resubmitted components (\S\ref{sec:resubmission}), an entire class of
Authenticode certificate-chain endpoints (\code{crl.comodoca.com},
\code{cacerts.digicert.com}) present in essentially any code-signed PE.
None of these are adversary-chosen infrastructure, yet every STIX bundle
exported before the fix promoted them to full \code{domain-name}
indicators regardless. A TLD allowlist, a compound-identifier casing
check (\code{MyApplication.app} is PascalCase in a way a human-typed C2
domain essentially never is), and a small, exact-matched allowlist for
the confirmed-benign PKI endpoints removed every instance across all
three families and all 26 resubmitted components.

Second, and unrelated to the first, CAPE's dynamic network capture
(\code{network.hosts}) is a whole-VM PCAP dump with no per-process
attribution: it cannot distinguish the sample's own connections from the
sandbox guest's background OS traffic. Cross-referencing every IP
observed across all three families' detonations against ASN registration
data found 18, none tied to any domain the sample itself resolved,
resolving to Microsoft Corporation (Azure AD/Teams telemetry) or Akamai
(the CDN fronting Windows Update and Defender cloud lookups), and
repeating byte-for-byte across a ransomware family, an infostealer, and
a loader, three samples with no plausible reason to share
infrastructure. A narrowly-scoped allowlist of these specific,
individually-confirmed IPs, deliberately not a blanket ASN or CIDR rule
(an Azure-Front-Door-fronted hostname elsewhere in the same dataset was
left untouched, since domain fronting through a major CDN is a real
adversary technique and ASN alone is not evidence of innocence there),
removed them from the reported IOCs.

The fix is independently checkable against WhiteSnakeStealer's own
ground truth, which documents a 31-node, 29-unique-IP HTTP C2 list
recovered by manual reverse engineering, the same protocol that
produced Table~\ref{tab:validation}. After the sandbox-noise filter,
the pipeline's dynamic-capture output recovers that list exactly: 29 of
29, zero missed. CAPE's own \code{dead\_connect} signature (confidence
100\%, ``31 unique'' unreachable attempts) independently corroborates
why every one of these reads as unreachable in the sandbox trace:
containment-first sandboxing (\S\ref{sec:limitations}, item~\ref{lim:containment}) means the
sample's own connection attempts to its real C2 fail by design, not
because the recovered addresses are fabricated. Contamination was not
fully zero, though: re-auditing this exact claim during the migration
to ATT\&CK v19 (2026-08) found 5 further IPs the two filters above did not
catch, none affecting the technique-level scores above since IOC
filtering and technique scoring are independent. Two, resolving to
self-identifying \code{*.msedge.net} hostnames on Microsoft's own ASN,
were confirmed and added to the allowlist the same way the existing
entries were. The remaining two, Cloudflare-range addresses with no
hostname, needed a stronger check than ASN and cross-sample repetition
alone before joining a list this project deliberately keeps to
zero-ambiguity entries: cross-referencing CAPE's raw per-process API
call trace, not just the summarized capture, confirmed neither IP
appears in the sample process's own 38 \code{connect()} calls (all 29
genuine C2 addresses, zero matches) nor in any other monitored
process's calls at all, only in the whole-VM packet capture -- no
hooked process is responsible for either connection, the same
sandbox-traffic signature the Microsoft entries were confirmed against.
Neither is a string in either binary or either manual report, ruling
out an as-yet-undecoded embedded destination. Both were then added. The
fifth, resolving to \code{ip-api.com}, is not contamination at all: the
same process-level trace shows \code{sample.exe} itself calling
\code{WinHttpGetProxyForUrl} against
\code{ip-api.com/line?fields=query,country}, a public IP-geolocation
lookup with no corresponding literal string anywhere in the binary, the
URL evidently built at runtime. This is a genuine, sample-initiated
finding, not a pipeline claim to trust at face value: added to
WhiteSnakeStealer's manual ground truth as
\technique{T1016} (System Network Configuration Discovery, whose own
MITRE page covers exactly this pattern, an online lookup for a
compromised system's own public-facing IP), the same evidentiary bar as
every other entry in that report, and the same way the \technique{T1055}
row earlier in the same report was already confirmed dynamically rather
than through static RE. Mapped to the conservative, literally-supported
technique rather than folded into the existing ``12 VM indicators''
anti-sandbox finding, since no conditional branch on the returned
country was independently observed to justify that stronger claim.

The lesson generalizes the same way the other cases do: every exclusion
is a specific, individually-justified, auditable entry, logged with its
own reason string, not a general reputation filter, since the whole
point of this pipeline over a reputation aggregator
(\S\ref{sec:related-work}) is that every decision stays traceable to a
concrete cause rather than asserted.

The fix's own generality was checked, not assumed: running the same
discipline against a fourth, structurally unrelated family (a
native-crypter-wrapped .NET RAT/backdoor, added specifically to test
generalization past the three validated ones) surfaced two further
instances of the same failure mode on that family's own resubmitted
components, a XOR-recovered value shaped like a .NET version tuple and a
second repeating-byte-padding cascade under a different key, plus a
class never seen before: X.509 certificate OIDs (\code{2.5.4.*},
\code{2.5.29.*}) mistaken for IPs on two of RoningLoader's own
resubmitted components, found by the same sweep. Each was fixed the same
way: a general, evidenced rule, verified against every known-real IP
across both families before being trusted.
